% META: {"id": "TEXP-MAT-ME-003", "chapitre": "TEXP-MATRICES-MARKOV", "type_objet": "methode", "capacites_codes": ["C3"], "methodes": ["M3"], "mode_creation": "ex_nihilo", "sources_inspiration": [], "status": "needs_review"}
% BEGIN-VERIFY
% from sympy import Matrix, Rational, eye
% A = Matrix([[Rational(1,2), 0], [0, Rational(1,3)]])
% C = Matrix([1, 2])
% L = (eye(2) - A).inv()*C
% assert L == Matrix([2, 3])
% U0 = Matrix([4, 6])
% U1 = A*U0 + C; U2 = A*U1 + C
% assert U1 == Matrix([3, 4]) and U2 == Matrix([Rational(5,2), Rational(10,3)])
% assert A**2*(U0 - L) + L == U2
% END-VERIFY
\begin{fichemethode}{M3}{Étudier une suite U(n+1) = A U(n) + C}

\quandUtiliser{%
  Utiliser cette méthode lorsqu'une suite de matrices colonnes vérifie
  $U_{n+1} = A\,U_n + C$ et que l'énoncé demande une expression de $U_n$, son
  comportement, ou un etat stable.
}

\pasApas{%
  \begin{enumerate}
    \item Chercher l'etat stable $L$ : résoudre $L = A\,L + C$, c'est-a-dire
      le système linéaire $(I - A)L = C$ (inverse de $I - A$ ou résolution
      directe).
    \item Poser la suite auxiliaire $V_n = U_n - L$.
    \item Montrer que $V_{n+1} = A\,V_n$ :
      $V_{n+1} = A U_n + C - (AL + C) = A(U_n - L)$.
    \item En déduire $V_n = A^n V_0$, puis
      $U_n = A^n\left(U_0 - L\right) + L$.
    \item Étudier le comportement : si $A^n \to 0$ (coefficients), alors
      $U_n \to L$ — interpréter l'etat stable dans le contexte.
  \end{enumerate}
}

\exempleRedige{%
  $U_{n+1} = A\,U_n + C$ avec
  $A = \begin{pmatrix} 1/2 & 0\\ 0 & 1/3 \end{pmatrix}$,
  $C = \begin{pmatrix} 1\\ 2 \end{pmatrix}$,
  $U_0 = \begin{pmatrix} 4\\ 6 \end{pmatrix}$. Exprimer $U_n$.

  \medskip
  \commentaireMarge{Etat stable : $(I - A)L = C$.}%
  $L = A L + C$ donne $\ell_1 = \frac{\ell_1}{2} + 1$ et
  $\ell_2 = \frac{\ell_2}{3} + 2$, soit $L = \begin{pmatrix} 2\\ 3
  \end{pmatrix}$.

  \medskip
  \commentaireMarge{Suite auxiliaire géométrique matricielle.}%
  $V_n = U_n - L$ vérifie $V_{n+1} = A V_n$, donc $V_n = A^n V_0$ avec
  $V_0 = \begin{pmatrix} 2\\ 3 \end{pmatrix}$ :
  $U_n = A^n \begin{pmatrix} 2\\ 3 \end{pmatrix} + \begin{pmatrix} 2\\ 3
  \end{pmatrix} = \begin{pmatrix} 2 + 2/2^n\\ 3 + 3/3^n \end{pmatrix}$.

  \medskip
  Les puissances $\frac{1}{2^n}$ et $\frac{1}{3^n}$ tendent vers 0 :
  $U_n \to L$.
}

\pieges{%
  \begin{itemize}
    \item \textbf{Piège 1.} Écrire $U_n = A^n U_0 + C$ : la constante ne se
      propage pas ainsi — passer par l'etat stable.
    \item \textbf{Piège 2.} Inverser $I - A$ sans vérifier qu'il est
      inversible ($\det(I - A) \neq 0$).
    \item \textbf{Piège 3.} Multiplier par $A$ du mauvais cote dans la suite
      auxiliaire ($V_{n+1} = A V_n$, l'ordre compte).
  \end{itemize}
}

\verifier{%
  Contrôler la formule sur $U_1$ et $U_2$ par la récurrence directe, vérifier
  $L = AL + C$, et confronter la limite au sens concret de l'etat stable.
}

\sEntrainer{\refExos{M3}}

\end{fichemethode}
